\documentclass[11pt]{article}

% ---------------------------------------------------------------------------
% The F2Z forest-GKR ALGORITHM note. Build: latexmk -pdf main.tex
% The protocol itself, in execution order, plus the lookup-table machinery
% the prover runs it on — deliberately free of scheduling/kernel detail.
% Companion notes: docs/forest-gkr-note/ (the prover engineering, measured),
% docs/verifier-note/ (every verifier check), docs/DESIGN.md (the pipeline).
% ---------------------------------------------------------------------------

\usepackage[margin=1.05in]{geometry}
\usepackage{amsmath,amssymb,amsthm}
\usepackage{booktabs}
\usepackage{array}
\usepackage{enumitem}
\usepackage{microtype}
\usepackage{xcolor}
\usepackage[colorlinks=true,linkcolor=blue!50!black,citecolor=blue!50!black,urlcolor=blue!50!black]{hyperref}

% --- notation (aligned with docs/forest-gkr-note/main.tex) ------------------
\newcommand{\F}{\mathbb{F}}
\newcommand{\Z}{\mathbb{Z}}
\newcommand{\K}{\mathbb{K}}            % the binding field GF(2^128)
\newcommand{\Kx}{\K^{\times}}
\newcommand{\Ftwo}{\F_2}
\newcommand{\Fq}{\F_q}                 % the evaluation field (char != 2)
\newcommand{\bits}{\{0,1\}}
\newcommand{\INT}{\mathsf{INT}}
\newcommand{\mle}[1]{\widetilde{#1}}
\newcommand{\eq}{\mathsf{eq}}
\newcommand{\cw}{c_w}
\newcommand{\absorb}{\textbf{absorb}}
\newcommand{\squeeze}{\textbf{squeeze}}
% \code: typewriter names; underscores allow a line break after them.
\newcommand{\code}[1]{{\def\_{\textunderscore\allowbreak}\texttt{#1}}}

\tolerance=1600
\setlength{\emergencystretch}{1.8em}
\hbadness=2500

\newtheorem{remark}{Remark}
\theoremstyle{definition}
\newtheorem{identity}{Identity}

\title{The F2Z forest GKR:\\ an algorithmic description}
\author{The \code{f2z-pcs} repository}
\date{2026-08-15}

\begin{document}
\maketitle

\begin{abstract}
\noindent
This is the \emph{algorithmic} companion to the prover-engineering note
\code{docs/forest-gkr-note}: the protocol the merged grand-product forest
performs, written in the order the prover and verifier execute it --- the
statement, the entry, the layer loop, the exit discharge into the
ring-switch/Ligerito opening, the read-off, and the soundness chain ---
followed by a self-contained account of the lookup-table machinery by
which the prover runs the bottom of the forest without ever materialising
it.  Kernel internals, software prefetch and the full measured cost tables
live in the companion note; nothing in the execution machinery described
here changes a transcript byte.  Everything here is implemented in
\code{src/merged\_forest.rs}, \code{src/piop/sumcheck/eq\_factored.rs}
and \code{src/ligerito.rs}; \S\ref{sec:code} maps the text to function
names.
\end{abstract}

\tableofcontents

% ===========================================================================
\section{What the forest certifies}
\label{sec:claim}

F2Z proves an integer multilinear-extension evaluation
$\mle{\INT(D)}(r) = y \in \Fq$ over data $D$ whose \emph{bits} are
committed with a characteristic-2 commitment.  The instance shape is
$(t, s, W)$: $D$ is a $2^t \times 2^s$ matrix of $W$-bit cells, arranged
as the single-bit matrix
\[
  M[c][i] \in \bits, \qquad
  i = (b \ll \log_2 W)\,|\,j \in \bits^{d}, \qquad
  d = t + \log_2 W,
\]
with $b$ the row, $c$ the column, $j$ the bit position inside the cell.
Splitting the evaluation point $r$ into its row and column halves, the
claim decomposes as
\[
  y \;=\; \sum_{c} \eq(c, r_{\mathrm{col}})
          \cdot u_c \bmod q,
  \qquad
  u_c \;=\; \sum_{b} w_b \cdot \INT\!\big(D(b,c)\big),
  \qquad
  w_b = \eq(b, r_{\mathrm{row}}) \bmod q
\]
with the row weights $w_b$ lifted to integers in $[0, q)$.  The prover
will simply \emph{send} the per-column integers and let the verifier do
the column combination in the clear; the entire difficulty is certifying
the sent integers against the $\Ftwo$ commitment, since parities destroy
magnitudes.  F2Z carries the integer fold in the exponent of
$\K = \mathrm{GF}(2^{128})$, with $\alpha \in \Kx$ a fixed public
generator.

For one chunk of the mod-$q$ opening (chunking is \S\ref{sec:chunks};
read this section at $L = 1$, i.e.\ $W_b = w_b$), with public $\cw$-bit
weights $W_b$, define the public \emph{leaf constants} and the
\emph{leaves}
\begin{equation}\label{eq:leaves}
  \tau(i) \;=\; \alpha^{W_b 2^j} - 1,
  \qquad
  \Lambda[c][i] \;=\; 1 + M[c][i]\cdot\tau(i)
  \;\in\; \{\,1,\ \alpha^{W_b 2^j}\,\}
\end{equation}
(in $\K$, $-1 = +1$; the $-$ form is kept for readability).  Then,
purely by construction,
\begin{equation}\label{eq:product}
  \prod_{i \in \bits^d} \Lambda[c][i]
  \;=\; \alpha^{\sum_{b,j} W_b 2^j\, M[c][i]}
  \;=\; \alpha^{u_c}.
\end{equation}
So ``the sent integers are the true weighted folds of the committed
bits'' \emph{is} the statement that $\alpha^{u_c}$ equals a grand
product of $2^d$ factors, each of which is $1$ or a public constant
\emph{selected by one committed bit} --- for all $2^s$ columns at once.
That is the claim the forest certifies.  The factors being
\emph{bit-affine} in the committed data (degree $1$ in each bit) is what
every later step leans on.

% ===========================================================================
\section{The protocol}
\label{sec:protocol}

The proof is a product-tree GKR in the lineage of
\cite{GKR15,Thaler13}: one binary product tree per column, walked from
the roots \emph{down} to the leaves, one sumcheck per tree level.  Three
F2Z-specific choices shape it:
(i) the roots are never sent --- the verifier recomputes them as
$\alpha^{u_c}$ from the sent integers;
(ii) the $2^s$ trees are \emph{merged}: the column index is just $s$
more multilinear variables, so one sumcheck per level covers every tree
and nothing per-tree ever enters the transcript;
(iii) the leaf layer is never committed --- because leaves are
bit-affine, the final leaf claim converts into a bare evaluation of the
committed bit matrix, discharged by the PCS opener.

\subsection{Entry: integers, roots, one claim}
\label{sec:entry}

\begin{enumerate}[leftmargin=2.2em, itemsep=2pt]
\item $P \to V$: the integer folds $u_c$ for every column $c$
  (16 bytes each; the \code{us} vector of the proof stream).
\item $V$: range-check every $u_c < 2^{\cw + t + W} = 2^{127}$, and
  check that $\alpha$ generates $\Kx$.  Together these make
  $x \mapsto \alpha^x$ \emph{injective} on every value that can occur
  ($\operatorname{ord} \alpha = 2^{128} - 1$).
\item Both sides: $R_c \leftarrow \alpha^{u_c}$ (the verifier by a
  fixed-base comb from the sent integers; the honest prover has them as
  the tops of its trees --- their equality is precisely the claim under
  proof), then
  \[
    \absorb(\mathtt{0x30}: R_0, \dots, R_{2^s-1}), \qquad
    \zeta \leftarrow \squeeze(\K^s), \qquad
    C \leftarrow \mle{R}(\zeta),
  \]
  each side computing the entry claim $C = \mle{R}(\zeta)$ itself
  ($O(2^s)$ work).  Set $z_x \leftarrow (\,)$, $z_c \leftarrow \zeta$.
\end{enumerate}

The absorb is the binding: a prover whose later messages were generated
against different roots faces challenges derived from \emph{these}
roots.  And instead of $2^s$ separate root claims, the verifier holds
from here on exactly \emph{one} running claim, about the MLE of the root
vector at a random point.

\subsection{The descent: one sumcheck per level}
\label{sec:descent}

Index the values of tree $c$ at depth $\ell$ (root $= 0$, leaves $= d$)
by $x \in \bits^\ell$, giving one multilinear family per level,
$L_\ell(x, c)$ with $x$ the low variables and $c$ the high ones, related
by the product recursion (children paired on the \emph{top} in-tree
bit: node $x$'s children sit at positions $x$ and $x + 2^\ell$ of level
$\ell+1$):
\begin{equation}\label{eq:recursion}
  L_\ell(x, c) \;=\;
  \underbrace{L_{\ell+1}(x, 0, c)}_{E_c(x)} \cdot
  \underbrace{L_{\ell+1}(x, 1, c)}_{O_c(x)}.
\end{equation}
$L_0(\cdot, c) = R_c$ are the roots and $L_d = \Lambda$ the leaves.

\medskip
\noindent
\textbf{Invariant.}  Entering layer $\ell$ the verifier holds one claim
$C = \mle{L_\ell}(z_x, z_c)$, $z_x \in \K^\ell$, $z_c \in \K^s$.  By
\eqref{eq:recursion} and multilinearity it factors coordinate-wise:
\begin{equation}\label{eq:layer}
  C \;=\; \sum_{x \in \bits^\ell} \sum_{c \in \bits^s}
     \eq(x, z_x)\,\eq(c, z_c)\; E_c(x)\, O_c(x).
\end{equation}
The layer converts this into the same kind of claim one level down:

\begin{enumerate}[leftmargin=2.2em, itemsep=2pt]
\item \textbf{Phase A} (the $\ell$ in-tree variables; skipped at
  $\ell = 0$).  A sumcheck over $x$ of \eqref{eq:layer}, binding
  $x \to r_x \in \K^\ell$; every round polynomial has degree 3
  ($\eq \cdot E \cdot O$).  Verifier checks: phase A's claimed sum
  equals the running claim $C$, and its final value equals
  $\eq(r_x, z_x) \cdot (\text{phase B's claimed sum})$.
\item \textbf{Phase B} (the $s$ tree variables).  A sumcheck over $c$
  of the residual $\sum_c \eq(c, z_c)\, \mle{E_c}(r_x)\,
  \mle{O_c}(r_x)$, binding $c \to r_c \in \K^s$.
\item \textbf{Closing pair and line step.}
  $P \to V$: the two children
  \[
    (p, q) \;=\; \big(\mle{L_{\ell+1}}(r_x, 0, r_c),\;
                      \mle{L_{\ell+1}}(r_x, 1, r_c)\big).
  \]
  Verifier checks phase B's final value $= \eq(r_c, z_c)\cdot p\cdot q$.
  Then both sides run
  \[
  \begin{gathered}
    \absorb(\mathtt{0x32}: p, q), \qquad
    \lambda \leftarrow \squeeze(\K), \\
    C \leftarrow p + \lambda\,(q - p), \qquad
    z_x \leftarrow r_x \| \lambda, \qquad
    z_c \leftarrow r_c .
  \end{gathered}
  \]
\end{enumerate}

The line step is the standard two-to-one reduction: the sumcheck leaves
\emph{two} evaluations of level $\ell+1$ (the children at
$(r_x, \cdot\,, r_c)$); folding them along the paired variable with a
fresh $\lambda$ gives back a single claim
$C = \mle{L_{\ell+1}}(r_x \| \lambda,\, r_c)$, restoring the invariant.
(The code calls $\lambda$ \code{mu}; $- = +$ in char 2.)

After $d$ layers the loop exits with the \emph{exit claim}
\[
  e_d \;=\; C \;\overset{\text{honest}}{=}\; \mle{\Lambda}(z_x, z_c),
  \qquad z = (z_x, z_c) \in \K^{d+s},
\]
one evaluation of the leaf MLE at a random point.

\begin{remark}[Merging]
A naive realisation runs $2^s$ independent GKRs and pays a per-tree
child pair every layer --- $O(2^s d)$ proof.  Treating the tree index as
$s$ more MLE variables makes the per-layer payload $O(\ell + s)$ round
messages plus \emph{one} child pair, and phase A's per-tree finals
$\big(\mle{E_c}(r_x), \mle{O_c}(r_x)\big)$ stay prover-internal ---
nothing per-tree is ever absorbed or sent.  At the deployed shapes the
sent folds ($16 \cdot 2^s$ bytes) are the larger forest-side proof item.
\end{remark}

\begin{remark}[The A/B split is bookkeeping, not protocol]
Semantically each layer is one $(\ell + s)$-variable sumcheck of
\eqref{eq:layer}; the implementation processes the $x$ variables first
(the $2^s$ trees run as eq-weighted groups sharing the point $z_x$, each
tree's $\eq(c, z_c)$ riding as a per-group scalar) and then the $c$
variables over the tiny per-tree finals.  The split is what the lazy
per-tree round bodies of \S\ref{sec:luts} and the across-tree
parallelism need; the transcript cost is one extra sumcheck header per
layer.
\end{remark}

\subsection{The exit claim and its discharge}
\label{sec:discharge}

The leaf layer never becomes an oracle.  Substituting
\eqref{eq:leaves} into the exit claim, splitting
$\eq((i,c), z) = \eq(i, z_x)\,\eq(c, z_c)$ and using
$\sum \eq = 1$:
\begin{equation}\label{eq:exit}
  e_d - 1
  \;=\; \sum_{i \in \bits^d} R[i]\cdot m[i],
  \qquad
  \begin{aligned}
    R[i] &= \eq(i, z_x)\,\tau(i)
      && \text{(public, size } 2^d\text{)},\\
    m[i] &= \textstyle\sum_c \eq(c, z_c)\, M[c][i]
      && \text{(the $z_c$-combined rows)}.
  \end{aligned}
\end{equation}
The exit claim is thus already a $\K$-linear functional of the committed
bits --- but against a \emph{dense, unstructured} weight vector $R$ (the
$\alpha$-powers destroy any tensor structure), which the opener cannot
consume: it proves plain MLE \emph{point} evaluations.  One last
\textbf{degree-2 sumcheck} (the ``presum'') separates the factors:

\begin{enumerate}[leftmargin=2.2em, itemsep=2pt]
\item Run a $d$-round degree-2 sumcheck on
  $\sum_i R(i)\, m(i) = e_d - 1$, binding $i \to r^*$.  (Verifier
  checks its claimed sum against the forest's $e_d$.)
\item $V$: evaluate $\mle{R}(r^*)$ \emph{itself} --- the verifier's one
  linear-cost step, $O(2^t \cdot W)$, with the $\alpha^{W_b}$ bases from
  the same fixed-base comb; reject if $\mle{R}(r^*) = 0$.  Set
  \[
    \mu \;\leftarrow\; \frac{\text{final sumcheck value}}{\mle{R}(r^*)}.
  \]
\item The residual claim is
  $\mle{M}(r^* \| z_c) = \mu$ --- a \emph{bare MLE evaluation of the
  committed bit matrix at a random point}, every $\alpha$-dependent
  factor moved to the verifier's side of the equation.
\end{enumerate}

From here the opening is out of the forest's hands: the ring switch
sends 128 partial evaluations $s_v$ (checked against $\mu$), the $L$
chunks' residual claims are $\eta$-combined, and \emph{one} recursive
Ligerito call certifies the batch against the Merkle commitment
(black box here; verifier note \S6).

\subsection{Read-off and the soundness chain}
\label{sec:soundness}

The verifier finishes in the clear, in $\Fq$:
\[
  y \;=\; \sum_c \eq(c, r_{\mathrm{col}}) \cdot
          \Big(\sum_l 2^{\cw l}\, u^{(l)}_c\Big) \bmod q,
\]
checked against the claimed value.  The chain, bottom-up --- each link
one of the steps above:
\begin{itemize}[itemsep=1pt]
\item \emph{Ligerito opening} $\Rightarrow$ $\mu$ is the true
  $\mle{M}(r^* \| z_c)$ of the committed bits;
\item \emph{presum $+$ the verifier's own $\mle{R}(r^*)$} $\Rightarrow$
  the exit claim $e_d$ is the true $\mle{\Lambda}(z)$;
\item \emph{$d$ layer sumchecks} $\Rightarrow$ $\mle{R}(\zeta)$ is the
  true MLE of the actual products $\prod_i \Lambda[c][i]$;
\item \emph{the absorb $+$ random $\zeta$} $\Rightarrow$ each root
  $R_c$ equals its product (Schwartz--Zippel over $\K$);
\item \emph{$R_c := \alpha^{u_c}$, the range check, the generator
  check} $\Rightarrow$ the sent $u_c$ \emph{are} the integer folds of
  the committed bits (injectivity of the exponent map below
  $\operatorname{ord}\alpha$);
\item \emph{recombination} $\Rightarrow$ $y = \mle{\INT(D)}(r) \bmod q$.
\end{itemize}
Both the range check and the generator check are load-bearing: without
either, $\alpha^{u} = \alpha^{u'}$ stops implying $u = u'$.

\medskip
\noindent
\textbf{Costs at a glance.}  Verifier: $O(2^s)$ exponentiations for the
roots $+$ $O(2^s)$ for the entry MLE and the read-off,
$O\big(\sum_\ell (\ell + s)\big) = O(d^2/2 + ds)$ round checks, and the
one $O(2^t W)$ evaluation of $\mle{R}(r^*)$ --- milliseconds at every
deployed shape.  Proof, forest side: the folds $16 \cdot 2^s$\,B plus
${\approx}\,32\big(\tfrac{d(d-1)}{2} + ds\big)$\,B of round messages and
$32d$\,B of pairs, per chunk.  Prover: the forest $+$ presum phase is
$O(2^{d+s})$ field work and 81--96\,\% of an F2Z prove; \S\ref{sec:luts}
is about how it is actually spent.

\subsection{Chunking: why there can be several forests}
\label{sec:chunks}

A full ${\sim}q_{\mathrm{bits}}$-bit weight $w_b$ would let
$\sum_b w_b\,\INT(D(b,c))$ overflow $2^{127}$ and break injectivity.  So
the weights are split into base-$2^{\cw}$ limbs, $\cw = 127 - t - W$,
$w_b = \sum_l 2^{\cw l}\, W^{(l)}_b$, and \emph{everything in
\S\S\ref{sec:entry}--\ref{sec:discharge} runs once per limb} on the same
transcript and the same commitment: per-limb integers $u^{(l)}_c$,
per-limb forest, per-limb presum, each chunk-fold satisfying
$u^{(l)}_c < 2^{\cw + t + W} = 2^{127}$ by construction.  The $L$
residual bit-MLE claims (at different points) are $\eta$-RLC'd into the
single Ligerito call, and the verifier reassembles
$\sum_l 2^{\cw l} u^{(l)}_c$ in the read-off.  At the deployed shapes
($q = 2^{100} - 15$, $t \le 26$) one has $\cw \ge 100$, so $L = 1$: the
whole scheme can be read with one chunk in mind.

% ===========================================================================
\section{Execution: the lookup-table prover}
\label{sec:luts}

Nothing below changes the protocol.  The tables exist to execute
\emph{phase A of the bottom three layers} (and the level build beneath
them) --- the rounds whose operand buffers $E_c, O_c$ would otherwise be
the two bottom levels of the forest, $2^{d+s}$ and $2^{d+s-1}$
$\K$-elements (4\,GiB $+$ 2\,GiB at $n = 28$).  Their messages and folds
are computed \emph{as if} those buffers existed, from lookups keyed by
committed bits; the transcript is byte-identical either way (pinned by
\code{lazy\_matches\_eager} in every schedule).

\subsection{What a round and a fold are}
\label{sec:roundfold}

Fix one layer $\ell$ and one tree, and look at its phase A through the
prover's eyes.  It is an ordinary $k$-variable sumcheck, $k = \ell$, of
\[
  S \;=\; \sum_{x \in \bits^k} w(x)\, E(x)\, O(x),
  \qquad w(x) = \eq(x, z_x),
\]
with $E, O$ --- the tree's two halves of the input level --- held as two
arrays of $2^k$ field elements.  The prover never manipulates a
polynomial symbolically; everything it does is one of two passes over
these arrays.

\paragraph{A round is one message--challenge exchange.}
The sumcheck runs $k$ rounds, one per variable, lowest index first
(a \emph{round} is a step \emph{within} one layer's phase A, not a
layer).  In round $j$ the prover sends the univariate
\[
  P_j(X) \;=\; \sum_{y \in \bits^{k-j}}
    \big(\mle{w}\,\mle{E}\,\mle{O}\big)
    (\rho_1, \dots, \rho_{j-1},\, X,\, y)
\]
--- the variables already bound are fixed at the earlier challenges, the
current one is left symbolic, the tail is still summed over the cube;
degree $3$ here, each factor being multilinear.  The verifier checks
$P_j(0) + P_j(1)$ against its running claim, draws $\rho_j$, and the
claim becomes $P_j(\rho_j)$.  After $k$ rounds nothing is left to sum:
the claim concerns the single random point
$r_x = (\rho_1, \dots, \rho_k)$.

\paragraph{The prover's state, and slots.}
$P_j$ is a sum, over the remaining cube, of products of
\emph{restrictions} of multilinears --- and a multilinear restricted to
$(\rho_1, \dots, \rho_{j-1}, X, b)$ is \emph{affine in $X$}, determined
by its two endpoints at $X = 0, 1$.  The prover therefore keeps each
array in the state ``bound prefix substituted, suffix still Boolean'',
\[
  E_j[y] \;=\; \mle{E}(\rho_1, \dots, \rho_{j-1},\, y),
  \qquad y \in \bits^{k-j+1}
\]
($E_1$ is the raw input), because those endpoints are then
\emph{literally stored entries}: grouping the entries into \emph{slots}
--- slot $b$ is the adjacent pair $(2b, 2b{+}1)$, the two settings of
the current variable ---
\[
  \mle{E}(\rho_1, \dots, \rho_{j-1},\, X,\, b)
  \;=\; E_j[2b] + X\,\big(E_j[2b{+}1] - E_j[2b]\big),
\]
and likewise for $O$ and $w$ --- though the $w$-line, unlike the other
two, will need no stored array at all (the closing paragraph of this
section) --- whence
\[
  P_j(X) \;=\; \sum_{b \in \bits^{k-j}}
    \big[\text{$w$-line}\big](X)\cdot
    \big[\text{$E$-line}\big](X)\cdot
    \big[\text{$O$-line}\big](X).
\]
One sweep over the slots --- two stored entries per slot per array ---
accumulates the message's coefficients: the $A_0, A_1, A_2$ of
\S\ref{sec:msgs}, the $w$ factor riding as the precomputed suffix
weights $\omega_b$.  (At Boolean points the entries \emph{are} the
current summand values, which is why $P_j(0) + P_j(1)$ matches the
running claim for the honest prover.)  Keeping the arrays in this state
is what makes the prover \emph{linear-time}: round $j$'s passes cost
$O(2^{k-j+1})$, the sizes halve, so a layer totals $O(2^k)$ ---
whereas recomputing the restrictions from the raw arrays each round
(each endpoint a fresh $2^{j-1}$-term $\eq$-weighted sum) would pay an
extra factor of $k$.

\paragraph{A fold is the array update after the challenge.}
Once $\rho_j$ arrives, substitute it --- evaluate every slot's line at
$X = \rho_j$:
\begin{equation}\label{eq:folddef}
  E_{j+1}[b] \;=\; E_j[2b] + \rho_j\,\big(E_j[2b{+}1] - E_j[2b]\big),
\end{equation}
one linear pass that merges each pair into one entry and \emph{halves}
the array, $2^{k-j+1} \to 2^{k-j}$.  Rounds are the protocol's clock;
folds are the bookkeeping tick that advances the arrays to the next
round's state.  A $k$-variable phase A thus walks
$2^k \to 2^{k-1} \to \dots \to 1$, and the last fold's single entries
are the per-tree finals $\mle{E}(r_x), \mle{O}(r_x)$ handed to phase B.
Unrolling \eqref{eq:folddef}: after $f$ folds,
\[
  E_{f+1}[b] \;=\; \sum_{\varepsilon \in \bits^{f}}
    \eq\big((\rho_1, \dots, \rho_f),\, \varepsilon\big)\cdot
    E_1[\,2^f b + \varepsilon\,],
\]
a \emph{fixed} $\K$-linear combination --- coefficients global across
trees and positions --- of one \emph{contiguous, aligned block} of $2^f$
round-1 entries.  (The alignment is why the LUT rounds below read their
case indices as aligned 2-bit, nibble and byte windows of the committed
bits.)

\paragraph{The $w$-line needs no table.}
The eq weight is the one factor whose line is never read from a stored
array, because $\eq$ factors coordinate-wise:
$\eq(x, z) = \prod_i \eq_1(x_i; z_i)$ with
$\eq_1(x; z) = zx + (1{-}z)(1{-}x)$, degree 1 in each variable
separately.  Writing $z_1, \dots, z_k$ for the coordinates of $z_x$,
the restriction to round $j$, slot $b$ splits along exactly those
lines:
\[
  \mle{w}(\rho_1, \dots, \rho_{j-1},\, X,\, b)
  \;=\;
  \underbrace{\prod_{i<j}\eq_1(\rho_i;\, z_i)}_{\textstyle A_j}
  \;\cdot\;
  \underbrace{\eq_1(X;\, z_j)}_{\text{public, slot-independent}}
  \;\cdot\;
  \underbrace{\prod_{i>j}\eq_1(b_{i-j};\, z_i)}_{\textstyle
    \omega_b = V_j[b]}
\]
and each piece has a different fate.  The \emph{prefix scalar} $A_j$
depends only on past challenges: one running field element, one
multiply per round ($A_{j+1} = A_j\,\eq_1(\rho_j; z_j)$) --- and in
the multi-tree phase A it is seeded with the tree's group scale
$\eq(c, z_c)$, which is how the merged layer's per-tree weights ride
along for free.  The \emph{current-variable factor} $\eq_1(X; z_j)$
carries all of the $w$-line's $X$-dependence and is the same public
line for every slot, so it is pulled outside the sum --- the prover
never multiplies by it at all (next paragraph).  The \emph{suffix
weight} $\omega_b = V_j[b]$ is the only per-slot part; since $b$ is
Boolean, each factor is $z_i$ or $1 - z_i$, so $V_j$ is a tensor, and
all $k$ of them are precomputed back-to-front ($V_k = [1]$,
$V_j = \eq_1(\cdot\,; z_{j+1}) \otimes V_{j+1}$;
\code{suffix\_tensors}) in $O(2^k)$ \emph{total} --- the cost of one
eq-table build --- once per phase-A call and shared by every tree,
since all trees share the point $z_x$.  The round polynomial
therefore factors,
\[
  P_j(X) \;=\; \eq_1(X;\, z_j)\cdot \hat H_j(X),
  \qquad
  \hat H_j \;=\; \sum_{c} A_{c,j} \sum_b \omega_b\,
    \big[\text{$E$-line}\big]\,\big[\text{$O$-line}\big],
\]
with $\hat H_j$ quadratic: the $(A_0, A_1, A_2)$ of \S\ref{sec:msgs}
are its per-tree coefficients, and the entire per-slot $w$ handling is
one scalar weighting folded into the $E$ side --- the 5-multiply
schedule $\omega e_0$, $\omega e_1$, $(\omega e_0)\,o_0$,
$(\omega e_1)\,o_1$, $\Delta E\,\Delta O$.  ``Folding $w$'' costs
nothing pass-shaped: one multiply into the prefix scalar and a step to
the next, half-length suffix tensor.

\paragraph{Where the public factor ends up: the verifier's side.}
Because $\eq_1(X; z_j)$ is public, the prover sends only
$\hat H_j$'s two non-constant coefficients $(\hat H_1, \hat H_2)$ ---
Gruen's factored round format.  The verifier recovers the constant
from its running claim: expanding
$S_j = P_j(0) + P_j(1) = (1{-}z_j)\,\hat H(0) + z_j\,\hat H(1)$, the
$\pm z_j \hat H_0$ terms cancel, leaving
\[
  \hat H_0 \;=\; S_j - z_j\,(\hat H_1 + \hat H_2)
\]
--- an identity over any field, no inversion --- after which it
advances $S_{j+1} = \eq_1(\rho_j;\, z_j)\cdot\hat H_j(\rho_j)$
(\code{verify\_eq\_inner\_sumcheck\_gruen}).  Soundness is unchanged:
the prover is thereby restricted to degree-3 polynomials divisible by
the public linear factor --- a \emph{subset} of the generic message
space, so the standard round analysis applies verbatim.

\subsection{Low entropy and the width laws}
\label{sec:width}

Everything rests on one observation: a node $r$ levels above the leaves
is the product of the $2^r$ leaves beneath it, and each leaf is $1$ or
the public constant $\alpha^{W_b 2^j}$ --- so the node's value is one of
$2^{2^r}$ \emph{position-indexed constants, selected by $2^r$ committed
bits, independent of the tree}.  One shared table can serve all $2^s$
trees; per tree one only reads bits.

The second ingredient is what a fold (\S\ref{sec:roundfold}) does to
such an entry: after $f$ folds, an array entry is a fixed $\K$-linear
combination of $2^f$ round-1 entries with globally shared coefficients
--- so if each round-1 entry was selected by $w$ committed bits out of a
shared position table, the folded entry is selected by the $2^f w$ bits
of its block, from a new position table with the challenges baked in.
Folding a bit-affine leaf pair makes this concrete --- using
$(1{+}\rho) + \rho = 1$ in char 2:
\begin{equation}\label{eq:fold}
  v' \;=\; (1{+}\rho)(1 + m_0\tau_0) + \rho\,(1 + m_1\tau_1)
      \;=\; 1 + m_0\,(1{+}\rho)\tau_0 + m_1\,\rho\,\tau_1 ,
\end{equation}
a 4-case function of the bit pair $(m_0, m_1)$, again from a
position-indexed, tree-shared table (now with $\rho$ baked in).  Hence
two \emph{width laws}, each of which squares the case count: going up a
level ($2^w \to 2^{2w}$ cases, the parent coupling two selects) and
folding a round (identically).  An entry of level $d - r$ after $f$
folds depends on $2^{r+f}$ committed bits --- $r$ \emph{fixed} for a
layer (how far above the leaves its input level sits), $f$
\emph{advancing} $0, 1, 2, \dots$ with the layer's own rounds --- and
its case table holds
$2^{2^{r+f}}$ cases per position:

\begin{center}
\small
\begin{tabular}{ccccl}
\toprule
$r + f$ & bits/entry & cases & table at $d = 17$ & \\
\midrule
0 & 1 & 2      & 2\,MiB  & the $\tau$ halves \\
1 & 2 & 4      & 4\,MiB  & \code{te}, \code{to} \\
2 & 4 & 16     & 8\,MiB  & $T4$ \\
3 & 8 & 256    & ${\approx}128$\,MiB & --- not built \\
4 & 16 & 65\,536 & ${\approx}8$\,GiB & --- exceeds what it would replace \\
\bottomrule
\end{tabular}
\end{center}

\paragraph{Why the cap is $r + f \le 2$.}
The cap is \emph{not} a soundness or correctness constraint --- any
depth would emit the identical transcript.  It is the economics knee of
a doubly-exponential blow-up ($+1$ \emph{squares} the case count), set
by three compounding facts.
(i) \emph{Cache residency.}  The value of a shared table is that all
$2^s$ trees gather from it with data-dependent (bit-selected) indices;
that pays only while the table is cache-resident.  The measured rule of
thumb is that one gather sweep over the 8\,MiB $T4$ already costs about
one \emph{streaming} pass over the level it replaces; a 128\,MiB table
turns every gather into a DRAM miss and costs \emph{more} than the dense
computation --- the rounds are gather-bound, not multiply-bound.  (Direct
evidence that the knee is residency: even within the cap, the 16-case
precombined leaf $\Delta\Delta$ table stops paying once it outgrows L2
and the code auto-switches to a factored 4-entry form, \S\ref{sec:msgs}.)
(ii) \emph{Build amortisation.}  A table costs one multiply per entry to
build; at $r{+}f = 3$ the build makes $32\times$ more values than the
level it serves has nodes per tree, eroding the $2^s$-fold amortisation
before the miss costs are even counted.
(iii) \emph{Diminishing returns.}  Each $+1$ buys one more round run off
bits --- one more halving of where dense buffers are born.  The 16-case
machinery already runs three rounds off the leaves; the marginal fourth
would halve an already-small object at the price of the 128\,MiB table.
Where a round \emph{message} couples several entries into a wider select
(the pair round's $\Delta E \cdot \Delta O$ term couples 8 bits), the
code does not build the 256-case table either: it \emph{factors} the
term into two 16-case tables joined by one explicit multiply.  Sixteen
cases is the last table size that earns its bytes; anything wider is
recovered by arithmetic on 16-case pieces.  (A separate, structural
guard --- a prefix of $f$ LUT rounds needs $k \ge f{+}1$ variables ---
just says the last LUT round must still have a fold in which the dense
buffers get born; it governs shallow trees, not the cap.)

\subsection{What the prover holds}
\label{sec:tables}

All shared tables are megabytes, amortised over the $2^s$ trees; the
per-tree state during the LUT rounds is \emph{only the bit arrays}.
Write $q_1 = 2^{d-2}$, $q_2 = 2^{d-1}$.

\begin{center}
\small
\begin{tabular}{l>{\raggedright\arraybackslash}p{5.6cm}lc}
\toprule
object & contents & index & at $d{=}17$ \\
\midrule
\code{lbits}/\code{rbits} (per tree) &
  the column's bits, split at the midpoint &
  bit $i$, 64/word & 32\,MiB total \\
$v$, $\tau$ halves &
  $v(i) = \alpha^{W_b 2^j}$, $\tau = v - 1$; case $0$ ($=1$) unstored &
  position $i$ & 2\,MiB \\
\code{pair} &
  $v(i)\,v(i{+}q_2)$ --- the only level-$(d{-}1)$ case not in $v$ &
  $i < q_2$ & 1\,MiB \\
\code{te}, \code{to} &
  the 4 cases $\{1, \tau_{\mathrm{lo}}, \tau_{\mathrm{hi}},
    \tau_{\mathrm{lo}}\tau_{\mathrm{hi}}\}$ of a level-$(d{-}1)$ node
  (E/O halves; O's offset baked in) &
  $4y \,|\, (m_{\mathrm{lo}} | m_{\mathrm{hi}}{\ll}1)$ & 4\,MiB \\
$T4$ &
  all 16 values of a level-$(d{-}2)$ node ($\code{te}\cdot\code{to}$) &
  $(y{\ll}4)|(c_E{\ll}2)|c_O$ & 8\,MiB \\
per-round tables &
  suffix-eq weights $\omega_b$ and past challenges baked into case
  tables (\S\ref{sec:msgs}--\ref{sec:stash}) &
  slot $\times$ case & MiB-scale \\
\bottomrule
\end{tabular}
\end{center}

The midpoint split of the bits is the top-bit pairing convention of
\eqref{eq:recursion} made physical: a level-$(d{-}1)$ node at position
$i$ multiplies leaves $i$ and $i + q_2$, and its two selecting bits are
bit $i$ of \code{lbits} and bit $i$ of \code{rbits} --- the same index
into both arrays, read as aligned bit windows word-wise.

\subsection{Use 1: round messages as gathers (the leaf round, worked)}
\label{sec:msgs}

A phase-A round accumulates three coefficients over slots $b$ (with
$\omega_b$ the driver's suffix-tensor eq weight --- $V_j[b]$ in the
code --- and entries $2b, 2b{+}1$ per side):
\[
  A_0 = \sum_b \omega_b\, E[2b]\,O[2b], \quad
  A_2 = \sum_b \omega_b\, \Delta E_b\, \Delta O_b, \quad
  A_1 = \Big(\sum_b \omega_b\, E[2b{+}1]\,O[2b{+}1]\Big) + A_0 + A_2
\]
(the $A_1$ form is the char-2 reading of $A_1 = H(1) - A_0 - A_2$).
At the leaf layer the entries are \emph{defined} as
$E[i] = 1 + l_i\,\tau_L[i]$, $O[i] = 1 + r_i\,\tau_R[i]$ ($l, r$ the
committed bit halves).  Expanding in char 2:
\begin{align*}
  \omega_b\,E[2b]\,O[2b]
  &\;=\; \omega_b \;+\;
     \big(\, l\cdot\omega_b\tau_L + r\cdot\omega_b\tau_R
        + lr\cdot\omega_b\tau_L\tau_R \,\big)
   && \text{a 4-case select,} \\
  \Delta E_b \;=\; E[2b{+}1] - E[2b]
  &\;=\; l_{2b}\,\tau_L[2b] + l_{2b+1}\,\tau_L[2b{+}1]
   && \text{the } 1\text{s \emph{cancel},} \\
  \omega_b\,\Delta E_b\,\Delta O_b
  &\;=\; \sum_{\varepsilon,\delta \in \bits}
     l_\varepsilon r_\delta\cdot
     \omega_b \tau_{L\varepsilon}\tau_{R\delta}
   && \text{a 16-case subset sum.}
\end{align*}
\paragraph{The three tables.}
Every slot contribution is a subset sum of \emph{eight} products ---
$\omega_b\tau_{L\varepsilon}$, $\omega_b\tau_{R\delta}$,
$\omega_b\tau_{L\varepsilon}\tau_{R\delta}$
($\varepsilon, \delta \in \bits$, all at slot $b$'s two positions) ---
plus a bare $\omega_b$ from each of the even and odd products.
\code{build\_leaf\_tables} precombines the subset sums, once per $\tau$
set, at 8 multiplies (4 singles, 4 crosses) and ${\sim}20$ adds per
slot:
\begin{itemize}[itemsep=2pt]
\item $\code{t\_a0}[(b{\ll}2) \,|\, (m_{L0} \,|\, m_{R0}{\ll}1)]$ ---
  the 4 cases
  $\{\,0,\ \omega_b\tau_{L0},\ \omega_b\tau_{R0},\
  \omega_b\tau_{L0} + \omega_b\tau_{R0} +
  \omega_b\tau_{L0}\tau_{R0}\,\}$
  of the even product's non-constant part;
\item \code{t\_a1} --- the same at the odd positions;
\item the $\Delta\Delta$ table,
  $\code{t\_a2}[(b{\ll}4) \,|\, c_l \,|\, (c_r{\ll}2)]$ --- the 16
  subset sums of the four cross products, selected by
  $m_{L\varepsilon} \wedge m_{R\delta}$;
\end{itemize}
and the constants collapse into one per-set scalar
$\code{w\_sum} = \sum_b \omega_b$ --- tree-independent, added once at
the end of the round.  Together: $4 + 4 + 16 = 24$ entries (384 B) per
slot, ${\approx}12.6$ MB at the $n = 28$ leaf layer ($k = 16$).

\paragraph{The slot body.}
A \emph{tree's} round is then, per slot:
\[
  \begin{aligned}
    &c_l \leftarrow \text{2-bit window of \code{lbits} at } 2b; \quad
     c_r \leftarrow \text{same in \code{rbits}}
     && \text{(32 slots per \code{u64} read)}\\
    &A_0 \mathrel{+}= \code{t\_a0}[(b{\ll}2)\,|\,c_{l,0}\,|\,c_{r,0}{\ll}1], \qquad
     S_1 \mathrel{+}= \code{t\_a1}[(b{\ll}2)\,|\,c_{l,1}\,|\,c_{r,1}{\ll}1], \\
    &A_2 \mathrel{+}= \Delta\Delta[(b{\ll}4)\,|\,c_l\,|\,c_r{\ll}2]
  \end{aligned}
\]
--- here $S_1$ collects the \emph{odd product sum}
$\sum_b \omega_b E[2b{+}1]O[2b{+}1] = H(1)$ (constants via
\code{w\_sum}), and the round ends with the finalisation
$A_1 = S_1 + A_0 + A_2$.  Per slot this is three unconditional indexed
loads and three XOR adds,
\textbf{zero multiplications, zero branches} on committed bits (a
branchy ``skip the multiply by $1$'' variant was measured slower:
50/50 random bits make it mispredict-bound).  The build is paid once
per set; the gathers are paid per tree: $O(2^k)$ shared multiplies
replace $O(2^k)$ multiplies \emph{per tree}, an amortisation of
$\mathit{live} \approx 2^s$.

\paragraph{Assembling the layer message.}
The triple a tree accumulates is its quadratic cofactor
$H_c = (A_0, A_1, A_2)$.  The driver multiplies each by the group's
running prefix scalar (seeded with $\eq(c, z_c)$,
\S\ref{sec:descent}) and XORs across trees into \emph{wide} unreduced
256-bit accumulators, with one polynomial reduction per accumulator per
round --- exact, because the reduction is $\Ftwo$-linear.  Since the
$\eq$ factor of the round polynomial is public, the message sent is
just the two non-constant coefficients of the summed cofactor
(Gruen's factored round format).

\paragraph{Two forms of the $\Delta\Delta$ table.}
The 16 $\Delta\Delta$ cases are subset sums of only four products, so
the table ships in two interchangeable forms (\code{LeafA2},
auto-picked by footprint; \code{F2Z\_LEAF\_A2\_FACTORED} forces):
\emph{precombined} --- 16 entries per slot, one load and one add ---
or \emph{factored} --- only the four raw cross products
$[p_{00}, p_{10}, p_{01}, p_{11}]$ per slot (one sequential 64-B
line), consumed by four branchless masked adds that reduce pairwise,
so the $A_2$ accumulator still takes one add.  Factored shrinks the
leaf tables from 24 to 12 entries per slot and wins outright once the
precombined set outgrows the P-cluster L2 (the auto threshold,
$2^{k-1} \ge 2^{15}$): at $n = 30$ the leaf round drops $2.1\times$
($307 \to 147$ ms); at the still-L2-resident $n = 26$ it loses
${\approx}1.2$ ms.  Both forms accumulate identical field elements.

\paragraph{One level up: the 16-case pair body.}
At the pair layer the entries are themselves 4-case
\code{te}/\code{to} selects, so a slot is keyed by four 2-bit cases
$(c_{E0}, c_{E1}, c_{O0}, c_{O1})$, read as 2-bit windows at $2b$ and
$2b + 2^k$ (\code{pair2\_cases}; the E/O offset is even, so no window
straddles a word).  The $\Delta E\cdot\Delta O$ term now couples
8 bits --- 256 cases, past the cap of \S\ref{sec:width} --- so it is
never fully precombined.  Two exact forms exist
(\code{Pair2Tables}):
\begin{itemize}[itemsep=2pt]
\item \emph{precombined} (opt-out, \code{F2Z\_PAIR2\_FACTORED=0}):
  four 16-case tables per slot ---
  $\code{t\_a0}[b{\ll}4 \,|\, (c_{E0}{\ll}2 \,|\, c_{O0})]$
  ${} = \omega_b\,TE_{2b}[c_{E0}]\,TO_{2b}[c_{O0}]$, the odd \code{t\_a1},
  the weighted difference $\code{t\_wde}$
  ${} = \omega_b\big(TE_{2b}[c_{E0}] + TE_{2b+1}[c_{E1}]\big)$, and
  the raw \code{t\_do} --- so the cross term
  is ONE wide multiply per slot; but the tables hold 64 entries per
  slot (16 MiB at the $n = 28$ pair layer, past L2), and the four
  gather streams' misses dominate the round;
\item \emph{factored} (the default): only the suffix-weighted
  $\code{wte}[i] = \omega_{i \gg 3}\cdot\code{te}[i]$ is built (the
  layout of \code{te} itself, 8 entries per slot), and the body
  recombines against the set's raw \code{to}: per slot, the weighted
  even and odd wide products plus
  $(\code{wte}_0 + \code{wte}_1)(\code{to}_0 + \code{to}_1)$ ---
  three wide multiplies against a $4\times$ smaller, L2-resident
  table set.
\end{itemize}
The two agree exactly: $\omega\,TE_0 + \omega\,TE_1$ \emph{is} the
precombined $\omega(TE_0 + TE_1)$ (distributivity), the products keep
the build's association $(\omega\,TE)\cdot TO$, and the deferred
reduction is $\Ftwo$-linear.  The measured lesson generalises across
this section: \emph{multiplies are cheap next to bytes that miss}.
This one body (\code{pair2\_factored\_body}) serves \emph{four}
rounds: round 1 of \code{Pair2Bits}/\code{Pair3Bits} over
$(\code{te}, \code{to})$, and round 2 of
\code{Leaf2Bits}/\code{Leaf3Bits} over the stashed value sets of
\S\ref{sec:stash} --- only the case extraction differs.

\subsection{Use 2: folds kept as tables (the stash)}
\label{sec:stash}

Round 1's fold should materialise round 2's buffers ---
$\mathit{live}\cdot 2^{k-1}$ $\K$-elements across the trees.  It does
not, because of one fact: \emph{the fold of a bit-selected buffer is
again a bit-selected buffer}.  So the prover never writes the folded
values --- it writes down the new (shared) selection table and keeps
using the bits.

\paragraph{Two representations of a buffer.}
A sumcheck operand can exist \emph{dense} --- an array of
$\K$-elements per tree, as \S\ref{sec:roundfold} assumed --- or in
\emph{case form}: one shared position-indexed table $T$ plus the
per-tree bits, under the convention that entry $p$ of tree $c$'s
buffer \emph{is defined as} $T[p][\,\text{a few of $c$'s bits at fixed
positions}\,]$.  In case form the buffer is a \emph{rule}, not memory:
no per-tree $\K$-element exists anywhere.  At the leaf layer the
buffers start in case form for free --- $E[i] = 1 + m_i\,\tau_L[i]$ is
the 2-case rule $\{1,\ 1{+}\tau_L[i]\}$ selected by one bit --- and
\S\ref{sec:msgs} already showed how a \emph{message pass} runs on that
form, by gathers.  What remains is the fold, and this is where the two
representations differ: folding a dense buffer is a write pass
\emph{per tree}; folding a rule is a rewrite of the rule --- one new
table, no tree touched.

\paragraph{Folding the rule.}
Substitute the case form into the fold \eqref{eq:folddef} at challenge
$\rho_1$ (char 2 cancels the $1$s, as in \eqref{eq:fold}):
\[
  E'[p] \;=\; (1{+}\rho_1)\,E[2p] + \rho_1\,E[2p{+}1]
  \;=\; 1 + m_{2p}\,(1{+}\rho_1)\tau_L[2p]
          + m_{2p+1}\,\rho_1\tau_L[2p{+}1].
\]
The right-hand side is a function of the bit \emph{pair}
$(m_{2p}, m_{2p+1})$, and its four possible values depend only on the
position and $\rho_1$ --- not on the tree.  So the fold of \emph{every
tree at once} is the one shared table
\[
  F_1[p][\,m_0 \,|\, m_1{\ll}1\,] \;=\;
  1 + m_0\,(1{+}\rho_1)\tau_L[2p] + m_1\,\rho_1\tau_L[2p{+}1]
\]
(4 entries per position, \code{build\_leaf\_fold\_tables}; one per
side, $F_1^L$ from $\tau_L$, $F_1^R$ from $\tau_R$), and
\emph{executing} it per tree is optional.  \code{LeafBits} does
execute it --- one $F_1$ gather per entry, writing dense round-2
buffers --- but \code{Leaf2Bits} and \code{Leaf3Bits} instead
\emph{stash} $F_1$: round 2's buffer of tree $c$ is henceforth
\emph{defined} as $p \mapsto F_1[p][m_{2p}, m_{2p+1}]$, and nothing
per-tree is written.  Round 2's message pass runs on the rule by
gathers: slot $b$ needs entries $2b, 2b{+}1$, i.e.\ the aligned nibble
$m_{4b..4b+3}$ of the \emph{original} bits, split into two 2-bit
cases (\code{leaf2\_cases}).  And by a layout coincidence the code
exploits, ``4 cases per position, two sides'' is exactly the pair
layer's native value format (\code{Pair2TauSet},
$\code{te} = F_1^L$, $\code{to} = F_1^R$) --- so round 2 \emph{is}
the pair-layer 16-case body of \S\ref{sec:msgs}, run over the stash.

\paragraph{Iterating, and stopping.}
At $\rho_2$ the same substitution folds the 4-case rule into a
16-case one:
\[
  E''[q] \;=\; (1{+}\rho_2)\,F_1[2q][c_0] + \rho_2\,F_1[2q{+}1][c_1]
  \;=:\; F_2[q][\,(c_0{\ll}2) \,|\, c_1\,],
\]
a function of the four original bits $m_{4q..4q+3}$
(\code{build\_pair2\_fold\_tables}, 16 entries per position).
\code{Leaf3Bits} stashes $F_2$ as round 3's value tables; round 3
gathers with byte windows ($m_{8b..8b+7}$ per slot, the nibble halves
swapping into the index via \code{leaf3\_idx}) and runs the plain
dense body over the gathered values --- at 8 selecting bits a
precombined slot table would cost more bytes than the multiplies it
saves (\S\ref{sec:width}).  There the chain stops: $\rho_3$'s fold
would need a 256-case $F_3$ --- each stash step \emph{squares} the
cases while only halving the positions, so the shared table
\emph{doubles} every step, the width law of \S\ref{sec:width} marching
toward its cap --- so round 3's fold is \emph{executed}: per tree,
gather the two $F_2$ values of each slot and write
$(1{+}\rho_3)(\cdot) + \rho_3(\cdot)$ into real arrays.  These are the
first per-tree $\K$-elements of the whole layer, $2^{k-3}$ per side;
rounds 4 onwards are ordinary dense sumcheck.
\code{Pair2Bits}/\code{Pair3Bits} are the same chain started one level
up, where the initial buffer is already a 2-bit-selected rule
$(\code{te}, \code{to})$, so the counting begins at 4 cases and the
cap arrives one round sooner (\code{pair3\_idx} keys the stashed
round-2 rule).

\paragraph{The prover's physical state, round by round.}
For \code{Leaf3Bits}, one $\tau$ set; table sizes count both sides, at
the $n = 28$ leaf layer ($k = 16$, $\mathit{live} = 2048$):

\begin{center}
\small
\begin{tabular}{lllr}
\toprule
 & round $j$'s buffer exists as & shared table & at $k{=}16$ \\
\midrule
round 1 & bits $+$ $\tau$ halves: $1 + m_i\,\tau[i]$ &
  $2\cdot 2^k$ entries & 2\,MiB \\
round 2 & bits $+$ $F_1[p][\,m_{2p}, m_{2p+1}\,]$ &
  $2\cdot 4\cdot 2^{k-1}$ & 4\,MiB \\
round 3 & bits $+$ $F_2[q][\,m_{4q..4q+3}\,]$ &
  $2\cdot 16\cdot 2^{k-2}$ & 8\,MiB \\
rounds 4+ & dense per-tree buffers & --- & 512\,MiB \\
\bottomrule
\end{tabular}
\end{center}

Until round 3's fold, the only per-tree state is the original bit
array; the challenges live inside the tables.  The cascade, per
driver entry (``stash'' = the round's fold is kept as the next round's
value tables; ``dense ($\cdot$)'' = the fold materialises buffers of
that size per side per tree):

\begin{center}
\small
\begin{tabular}{lllll}
\toprule
entry & needs & round 1 & round 2 & round 3 \\
\midrule
\code{LeafBits}  & $k \ge 2$ & leaf LUT & dense ($2^{k-1}$) & \\
\code{Leaf2Bits} & $k \ge 3$ & leaf LUT, stash & pair body &
  dense ($2^{k-2}$) \\
\code{Leaf3Bits} & $k \ge 4$ & leaf LUT, stash & pair body, stash &
  inline; dense ($2^{k-3}$) \\
\code{Pair2Bits} & $k \ge 2$ & pair body & dense ($2^{k-1}$) & \\
\code{Pair3Bits} & $k \ge 3$ & pair body, stash &
  inline; dense ($2^{k-2}$) & \\
\code{T4Bits}    & $k \ge 2$ & inline $T4$; dense ($2^{k-1}$) & & \\
\bottomrule
\end{tabular}
\end{center}

The ``needs'' column is \S\ref{sec:width}'s structural guard ($k \ge
f{+}1$: the last LUT round must still have a fold, where the dense
buffers get born).  \code{F2Z\_LUT3=0} forces the shallower
\code{Pair2}/\code{Leaf2} prefixes at every depth --- the same path
depth-4 trees take; depth-3 trees run \code{LeafBits} alone.

\paragraph{Reading the cases.}
Each round down a cascade reads a wider aligned bit window per slot,
always keyed by the \emph{original} committed bits --- nothing
per-tree is ever written between rounds, which is the point:

\begin{center}
\small
\begin{tabular}{lll}
\toprule
round body & windows per slot & helper \\
\midrule
leaf round 1 & two 2-bit fields at $2b$ & inline \\
\code{Leaf2}/\code{Leaf3} round 2 & two nibbles at $4b$ &
  \code{leaf2\_cases} \\
\code{Leaf3} round 3 & two bytes at $8b$ & \code{leaf3\_idx} \\
\code{Pair2}/\code{Pair3} round 1 &
  four 2-bit fields at $2b,\ 2b{+}2^k$ & \code{pair2\_cases} \\
\code{Pair3} round 2 & four nibbles at $4b,\ 4b{+}2^k$ &
  \code{pair3\_idx} \\
\code{T4Bits} round 1 & 16 single bits & \code{t4bits\_idx} \\
\bottomrule
\end{tabular}
\end{center}

The windows are \S\ref{sec:roundfold}'s contiguous-block fact made
physical: after $f$ folds an entry's selecting bits are one aligned
$2^{f}{\cdot}w$-bit block, and the E/O offset $2^k$ stays even at
every round, so no window ever straddles a word.

\paragraph{Economics.}
The count that decides everything: an executed fold stores one value
per position \emph{per tree}; a stashed one stores 4 (resp.\ 16)
values per position \emph{in total}.  Materialising at round $j$ thus
costs $\mathit{live}\cdot 2^{k-j}$ values per side; the stash costs
$4\cdot 2^{k-j}$ resp.\ $16\cdot 2^{k-j}$
entries per side \emph{per $\tau$ set} --- factors $\mathit{live}/4$
and $\mathit{live}/16$, i.e.\ $512\times$ and $128\times$ at the
deployed $s = 11$.  Concretely, at the $n = 28$ leaf layer ($k = 16$,
$\mathit{live} = 2048$): executing round 1's fold would write 2 GiB of
round-2 buffers, where the stash is 4 MiB; round 2's stash is 8 MiB
against a 1 GiB materialisation; and the first dense buffers, at
round 3's fold, are $2^{n-3}\cdot 16$ B $= 512$ MiB.  It is this LUT
prefix depth, not the storage schedule, that fixes the bottom layers'
resident set.

\subsection{Use 3: generating the dense region from $T4$}
\label{sec:jit}

Above the pair layer the entropy argument runs out ($r + f \le 2$), so
$T4$ is used not to \emph{avoid} buffers but to \emph{create} them
cheaply --- and the remaining freedom, which level is stored, which
regenerated, which implicit, is the memory schedule.

\paragraph{The stored chain, and the build.}
\code{build\_levels} builds each tree's upper levels bottom-up by
repeated \code{parent\_halves\_top} ($v[i] = E[i]\cdot O[i]$,
re-split at the new midpoint), parallel across trees, then regroups the
chains level-major; the roots fall out of level 1.  The chain's
\emph{top} level is produced directly from bits and $T4$
(\code{gen\_top}), one subcube of the level-width observation
(\S\ref{sec:width}) per position: under L/2 one $T4$ gather ($=$ level
$d{-}2$); under L/4 the product of the pair $\{y,\ y{+}2^{d-3}\}$ ---
2 gathers, 1 multiply ($=$ level $d{-}3$); under L/8 the quadruple
$\{y,\ y{+}2^{d-4},\ y{+}2^{d-3},\ y{+}3\cdot 2^{d-4}\}$ --- 4 gathers,
3 multiplies ($=$ level $d{-}4$).  Levels $d$ and $d{-}1$ (and, below
L/2, $d{-}2$) therefore never exist at build time: each of their
positions is gathered exactly once, with the materialise-then-read
round trip removed.  The chain that remains sums to
${\approx}\,2^{\mathrm{top}+1}$ values per tree --- a quarter of the
leaf count under L/4, an eighth under L/8, whence the names --- and
\code{drive\_grouped} takes each level with \code{mem::take} as its
layer starts, so the chain shrinks as the forest climbs.  The stored
chain and the JIT level below it are equal-sized and never
simultaneously live.

\paragraph{The JIT level, and fused generation.}
Under L/4, layer $d{-}3$ regenerates its input level $d{-}2$ per tree
($2^{d-3}$ values per half, one 4-bit-indexed $T4$ gather each) into
dense buffers that die when the layer ends; under L/8, layer $d{-}4$
does the same for level $d{-}3$ at two gathers and a multiply per
value.  By default the generation pass is not just a write:
\code{jit\_layer\_generate} accumulates the layer's opening round(s)
\emph{as it writes}.  In the plain form
(\code{dense\_jit\_fused\_round1}) that is round 1's coefficient
triple; with the driver's double-fold on and $k \ge 3$
(\code{dense\_jit\_fused\_grid}) it is the whole $3{\times}3$
bivariate grid
$G(X_1, X_2) = \sum_b V(b)\,L(X_1, X_2, b)\,R(X_1, X_2, b)$ --- nine
field elements per tree, from which round 1's message
$\sum_{x_2} \eq_1(x_2;\, q)\,G(X_1, x_2)$ \emph{and} round 2's
$G(\rho_1, X_2)$ both derive, so the level's first actual read is
round 3's pass, which then folds two deferred challenges at once.
Either way the result reaches the driver through its
\code{pre\_round1} hook and is value-exact against the passes it
replaces: XOR accumulation is order-free and the deferred reduction
$\Ftwo$-linear.

\paragraph{The 4-leaf round (L/8 only).}
One level further up, \code{T4Bits} runs a whole layer off bits:
each entry is a single inline 16-case $T4$ gather ---
\code{t4bits\_idx} assembles $(j{\ll}4) \,|\, (c_E{\ll}2) \,|\, c_O$
from four committed bits read straight from the transposed leaf
halves at $j$ and $j + 2^{k+1}$ --- with dense buffers at its fold.
The layer's input level is neither stored nor regenerated; combined
with the JIT level below it, this is what makes L/8 an eighth.

\paragraph{The three schedules.}
The schedules differ in nothing but the fate of levels $d{-}2$ and
$d{-}3$:

\begin{center}
\small
\begin{tabular}{llll}
\toprule
 & L/2 (\code{l2}) & L/4 (default) & L/8 (\code{l8}) \\
\midrule
chain tops at & level $d{-}2$ & level $d{-}3$ & level $d{-}4$ \\
per top position & 1 $T4$ gather & 2 gathers, 1 mul & 4 gathers, 3 muls \\
level $d{-}2$ & stored & JIT at layer $d{-}3$ & implicit (\code{T4Bits}) \\
level $d{-}3$ & stored & stored & JIT at layer $d{-}4$ \\
peak resident & ${\approx}2^n{\cdot}8$\,B & ${\approx}2^n{\cdot}4$\,B &
  ${\approx}2^n{\cdot}2$\,B \\
needs depth & $\ge 4$ & $\ge 4$ & $\ge 5$ (else L/4) \\
\bottomrule
\end{tabular}
\end{center}

Schedule choice is a near-pure \emph{memory} decision.  The measured
rule of thumb: one $T4$ gather sweep costs about one streaming pass
over the level it would have saved --- so L/2, which removes the JIT
sweep but doubles the resident set, comes out flat to worse than L/4,
while L/8 costs $+5$--$13\,\%$ of prove and is what makes $n \ge 31$
fit a 16\,GB box.  The time levers live inside the round bodies of
\S\S\ref{sec:msgs}--\ref{sec:stash}, not in the ladder.

\paragraph{Only live columns are built.}
A committed column whose bits are all zero has every leaf
$\alpha^0 = 1$, hence every level and its root $= 1$.  Since a phase-A
round message is $\sum_c \eq(c, z_c)\,H_c$ and every constant-1 tree
has the \emph{same} $H_c$, the trailing run of zero columns collapses
into \emph{one} synthetic group with all-ones buffers and scale
$\sum_{c \ge \mathit{live}} \eq(c, z_c)$ (\code{col\_elide},
\code{live\_cols}).  The absorbed root vector is still padded with the
known $1$s, so the statement is unchanged; on a JIT layer the
synthetic group is born inside the generation pass.  Char-2
re-association is exact, so the transcript is byte-identical to the
un-elided forest (\code{col\_elision\_matches\_full}): a witness of
$N$ cells padded to $2^n$ pays the forest for
$\lceil N / 2^{t + \log_2 W} \rceil$ columns instead of $2^s$.
Everything in this section runs over the live leading trees only.

\paragraph{The ladder, per layer.}
Seen from the layers, the default schedule reads:

\begin{center}
\small
\begin{tabular}{lllll}
\toprule
layer & input level & entries exist as & LUT rounds & dense born \\
\midrule
$d{-}1$ (leaf) & leaves & $1 + m\tau$ over the $\tau$ halves &
  1--3 (2-bit, nibble, byte) & round 3's fold \\
$d{-}2$ (pair) & $d{-}1$ & 4-case \code{te}/\code{to} selects &
  1, 2 & round 2's fold \\
$d{-}3$ & $d{-}2$ & JIT from $T4$ gathers & (fused into generation) &
  at generation \\
$\le d{-}4$ & $\le d{-}3$ & stored chain & --- & prebuilt \\
\bottomrule
\end{tabular}
\end{center}

\subsection{Why the transcript is unchanged}
\label{sec:identity}

Every device above is a re-association of the same field arithmetic:
precombined entries are subset sums the dense path computes per entry
(distributivity); the char-2 identities $1 + 1 = 0$ and
$(1{+}\rho) + \rho = 1$ are what cancel the leaves' $1$s and keep folded
entries affine; XOR accumulation is order-free and the deferred wide
reduction $\Ftwo$-linear.  Same field elements, same transcript order;
the verifier and the proof codec never learn which path ran.  The
byte-identity is pinned in every schedule and depth regime
(\code{lazy\_matches\_eager}, \code{lazy\_multi\_matches\_eager},
\code{col\_elision\_matches\_full}).

% ===========================================================================
\section{Where this sits in the code}
\label{sec:code}

\begin{center}
\small
\begin{tabular}{>{\raggedright\arraybackslash}p{5.9cm}>{\raggedright\arraybackslash}p{8.9cm}}
\toprule
this note & code \\
\midrule
the layer loop (\S\ref{sec:descent}) &
  \code{drive\_grouped} / \code{verify\_merged\_forest} \\
per-layer proof (phase A, phase B, $(p,q)$) &
  \code{MergedLayer} (\code{sc\_x}, \code{sc\_c}, \code{pair}) \\
the line challenge $\lambda$ & \code{mu} \\
the integer folds $u_c$ & \code{fold\_values\_bits}; proof field \code{us} \\
$\alpha$-power table $v(i)$ & \code{chunk\_pow2\_table} \\
$R$ table / $m$ table (\S\ref{sec:discharge}) &
  \code{row\_bit\_weights} / \code{xi\_combined\_rows} \\
the presum & \code{MultiDegreeSumcheck} in
  \code{prove/verify\_int\_eval\_merged\_common} \\
range/generator/root checks, $\mu$ extraction &
  \code{verify\_int\_eval\_merged\_common} \\
chunk loop, $\eta$-RLC, the one Ligerito call &
  \code{prove/verify\_mle\_eval\_mod\_q\_ligerito} \\
the lazy prover & \code{prove\_merged\_forest\_lazy} \\
LUT round bodies (\S\ref{sec:msgs}--\ref{sec:stash}) &
  \code{GroupBufs}: \code{LeafBits}, \code{Leaf2Bits}, \code{Leaf3Bits},
  \code{Pair2Bits}, \code{Pair3Bits}, \code{T4Bits} \\
leaf tables $\code{t\_a0}/\code{t\_a1}/\Delta\Delta/\code{w\_sum}$ &
  \code{build\_leaf\_tables} \\
fold/stash tables & \code{build\_leaf\_fold\_tables},
  \code{build\_pair2\_fold\_tables} \\
shared level tables & \code{leaf\_tau\_halves},
  \code{layer1\_pair\_table}, \code{te}/\code{to}, $T4$ \\
JIT generation (fused round 1 / grid) & \code{jit\_layer\_generate} \\
\bottomrule
\end{tabular}
\end{center}

For the checks in full detail see \code{docs/verifier-note/}; for the
schedules, kernels, fusion, and the measured cost tables see
\code{docs/forest-gkr-note/}; for the surrounding pipeline (commit,
ring switch, Ligerito, serialization) see \code{docs/DESIGN.md}.  The
standard caveats carry over: flock's Merkle domain separation is a
flagged upstream item, the ad-hoc Ligerito config regime is unaudited,
and the scheme is not zero-knowledge.

% ===========================================================================
\begin{thebibliography}{9}\small

\bibitem{GKR15} S.~Goldwasser, Y.~T.~Kalai, G.~N.~Rothblum.
\emph{Delegating computation: interactive proofs for muggles.}
J.~ACM 62(4), 2015 (STOC 2008).

\bibitem{Thaler13} J.~Thaler.
\emph{Time-optimal interactive proofs for circuit evaluation.}
CRYPTO 2013.  (The product-tree GKR specialisation the layers
instantiate.)

\bibitem{Libra} T.~Xie, J.~Zhang, Y.~Zhang, C.~Papamanthou, D.~Song.
\emph{Libra: succinct zero-knowledge proofs with optimal prover
computation.}  CRYPTO 2019.  (The eq-factored linear-time prover
lineage.)

\bibitem{Soukhanov} L.~Soukhanov.
\emph{The characteristic-2 field-switching note}, \S9.  (The
integer-MLE-over-$\Ftwo$ construction F2Z realises; internal.)

\end{thebibliography}

\end{document}
